\section{Summary of Contributions and Key Takeaways}
\label{sec:discussion:contributions}

\subsection{Poking as a Complementary Dynamic Manipulation Primitive}
\label{sec:discussion:poking}

\cref{ch:poking} introduced poking as a dynamic, non-prehensile manipulation primitive for object rearrangement in cluttered environments.
Where pushing relies on sustained contact and therefore couples object displacement to end-effector motion, poking uses brief impulsive contacts to transfer momentum to the object.
This distinction changes the geometry of the planning problem: a short robot motion can induce a comparatively large object motion, allowing the planner to reach object configurations that are inaccessible to pushing or grasping alone.

\textit{PokeRRT} operationalizes this idea through a closed-loop sampling-based kinodynamic planner that uses physics simulation as a forward model.
Rather than requiring an analytical model of impulsive contact, the planner samples candidate poking actions, propagates their effects through simulation, and replans as the observed object state evolves.
Across six experimental scenarios, this approach achieved higher success rates and shorter task completion times than pushing and grasping baselines, particularly in settings where grasping was geometrically infeasible or sustained-contact pushing was constrained by clutter and reachability.

The central contribution of this work is not only the poking primitive itself, but the demonstration that simulation-in-the-loop kinodynamic planning can expose useful manipulation behaviors that are difficult to capture with purely geometric abstractions.
Poking illustrates the broader premise of the thesis: dynamic interaction expands the set of feasible manipulation strategies when the planner is allowed to reason about state transitions, rather than only static configurations.

% \subsection{Poking as a Complementary Dynamic Manipulation Primitive}
% \label{sec:discussion:poking}

% Non-prehensile manipulation research before \cref{ch:poking} had focused almost entirely on pushing---contact modes in which the robot maintains continuous contact with the object throughout the trajectory.
% Pushing is limited because it ties object displacement to robot end-effector arc length: moving an object far requires the robot arm to travel far, and reachability constraints clip the feasible workspace accordingly.
% \textit{PokeRRT} introduced \textit{poking}---impulsive, instantaneous contact---as a complementary primitive, embedded in a closed-loop RRT-based planner that uses physics simulation as a forward model.
% The key insight is geometric: an impulsive contact interaction can displace an object substantially with a comparatively short end-effector trajectory, decoupling object travel distance from robot joint motion.
% In experiments across six scenarios, poking-based planners outperformed both pushing and grasping in success rate and achieved significantly lower task times by avoiding the constant-contact constraint.
% Poking succeeded in scenarios where grasping was geometrically infeasible and where pushing's monotonic displacement model caused failure in cluttered environments.
% The broader takeaway is methodological: simulation-in-the-loop planning can realize a manipulation skill whose contact dynamics would be prohibitively complex to model analytically, without requiring explicit contact models at planning time.

\subsection{Kinodynamic Planning for Whole-Body Contact and Coupled Dynamics}
\label{sec:discussion:constrained}

\cref{ch:constrained} extended simulation-in-the-loop kinodynamic planning from impulsive contact to manipulation problems in which the robot, object, and environment remain dynamically coupled throughout execution.
The chapter studied two such settings: manipulation under locked multi-joint failures, where the robot must exploit whole-body contact despite losing nominal end-effector functionality, and spill-free liquid transport, where feasible motion is constrained by motion-induced fluid dynamics.

In the locked-joint setting, the central insight is that actuation failure need not eliminate manipulation capability.
Rather than treating a failed robot as an exceptional case that must freeze or recover before continuing, the method treats each failure configuration as a new embodiment with its own feasible interaction modes.
By planning over whole-body non-prehensile contacts, the robot can repurpose links that are not normally considered manipulation surfaces and recover substantial workspace functionality even when the gripper is unavailable.
This reframes failure-aware manipulation as a kinodynamic planning problem over altered embodiment dynamics, rather than as a fault-handling problem outside the planner.

In the liquid-transport setting, the challenge is different but structurally related: feasible motion is constrained not only by collision avoidance and robot kinematics, but also by the coupled dynamics of the carried liquid.
The proposed planner uses a learned spill-free classifier as an implicit dynamics model, allowing it to reason about liquid safety without relying on simplified analytical approximations or direct surface monitoring.
This enables full six-degree-of-freedom container motion in cluttered environments while maintaining spill-free execution.

Together, these systems show that kinodynamic planning can be extended beyond isolated contact events to tasks where dynamic constraints persist throughout the motion.
They also demonstrate an important design principle: when explicit dynamics models are difficult to derive, useful models of dynamic feasibility can be supplied either by simulation, precomputed motion structure, or learned classifiers.
In this sense, \cref{ch:constrained} broadens the thesis argument from exploiting individual dynamic interactions to planning under sustained coupled dynamics.

% \subsection{Kinodynamic Planning for Whole-Body Contact and Coupled Dynamics}
% \label{sec:discussion:constrained}

% \cref{ch:constrained} extended kinodynamic planning to two settings where robot and object dynamics are coupled throughout the motion, not only at an impact instant.

% The first is manipulation under locked multi-joint failures (LMJs).
% Current systems respond to actuation failures by freezing---a conservative strategy that prioritizes safety but eliminates manipulation capability entirely.
% This work demonstrated that robots in failure states still possess substantial manipulation capability through whole-body non-prehensile contact.
% By pre-computing kinodynamic edge bundles for specific failure configurations and using a simulation-in-the-loop planner to select among them, the approach maintains up to 92\% of the nominal reachable area even when the gripper is entirely unusable, achieving a 79\% increase in failure-constrained reachable area and task success rates of 88.9--100\% across tested scenarios.
% The key insight is that failure modes are better reframed as new embodiment configurations to plan around than as exceptional states to recover from.

% The second contribution is spill-free liquid transport in cluttered environments (\textit{SFRRT*}).
% Prior approaches restricted container orientation to conservative angle ranges or used simplified linear liquid models (pendulum, mass-spring-damper) that cannot operate in genuinely cluttered scenes.
% \textit{SFRRT*} instead uses a transformer-based spill-free classifier trained on human-collected trajectory data as an implicit liquid dynamics model, combined with container geometry to estimate maximum safe tilt angles.
% This achieves a 3$\times$ expansion of the feasible solution space compared to prior approaches, with 93\%+ task success and full six-degree-of-freedom end-effector orientation flexibility---without explicit liquid modeling or visual surface monitoring.
% The takeaway is that implicit learned dynamics can substitute for explicit physics models in a safety-critical planning context, enabling capabilities that analytical models preclude.